\section{Introducción}
    Cuando se captura imagen digital de un objeto, no solamente se encuentra 
    sujeta a los factores de iluminación o colores para determinar su calidad 
    final sino también a las interferencias por fenómenos electromagnéticos a 
    momento de generar o transmitir la imagen digital, estas interferencias, 
    en la mayoría de los casos, se pueden observar como ruido periódico. Para
    tratar con este tipo de ruido se emplea de la transformada de Fourier, la 
    cual convierte una imagen (espacio de señal) a su espectro de frecuencias 
    o, equivalentemente, descompone una imagen en una combinación lineal de 
    senos y cosenos; y este es el punto clave para mitigar el ruido periódico, 
    como es una combinación de funciones senoidales se podrá detectar las 
    frecuencias asociadas con el ruido y eliminar (filtrarlas).\\

    Este proyecto tiene el fin de aplicar el filtrado de frecuencias asociadas 
    con ruido periódico en imágenes para mejorar su estética o, por al menos, 
    reducir el efecto que provoca en la imagen. Como se mostrará en los 
    resultados experimentales, emplear el espectro de frecuencias será una 
    forma eficiente para eliminar el ruido presente en una imagen sin tener 
    que elaborar convoluciones u otras técnicas más complejas.